\documentclass[10pt,a4paper]{article}

\usepackage[T1]{fontenc}
\usepackage[utf8]{inputenc}
\usepackage[margin=2cm]{geometry}
\usepackage{microtype}
\usepackage{xcolor}
\usepackage{hyperref}

\hypersetup{
  colorlinks=true,
  linkcolor=blue,
  filecolor=magenta,
  urlcolor=cyan,
  citecolor=blue,
  pdftitle={Proving Ethereum with Lattices (Extended Abstract)}
}

\setlength{\parindent}{1em}
\setlength{\parskip}{0pt}

\begin{document}
\begin{center}
	{\Large\bfseries Proving Ethereum with Lattices (Extended Abstract)\par}
	\vspace{0.35em}
	{\normalsize Luk\'{a}\v{s} \v{C}astven\par}
	{\small Slovak University of Technology in Bratislava\par}
\end{center}
\small
\thispagestyle{plain}

\noindent\textbf{Abstract.} Ethereum scalability is bottlenecked by the need
to repeatedly execute EVM computation on many nodes. Snarkifying execution
would replace much of that work with verification of a compact proof. Most
current ZkVMs obtain post-quantum security from hash based STARKs, while some
recursive systems still rely on elliptic curve assumptions. This work
studies a third path: a general purpose RISC-V ZkVM built from lattice
cryptography. The design combines CCS arithmetization, LatticeFold for
incrementally verifiable computation, deferred memory checking, and a
Greyhound wrapper proof. The evaluation is based on an incomplete Rust
implementation. The prototype is much slower than mature ZkVMs, but the design
is feasible, longer runs keep memory bounded, and the final proof in the design
is expected to be far smaller than current ZkVM proofs, in the tens of
kilobytes rather than hundreds of kilobytes to a megabyte.

\section*{Motivation}

Ethereum's long-term scalability depends on replacing widespread reexecution
with succinct proofs of correct computation \cite{buterin}. For execution,
this appears either as zkEVMs that model the EVM directly or as ZkVMs that
prove an easier virtual machine and embed Ethereum logic inside it
\cite{sokzkvms}. In practice, many systems have converged on RISC-V ZkVMs
because the EVM's design (stack oriented, use of Keccak256, \ldots)
is difficult to arithmetize efficiently.

At the proving layer, the design space is split between hash-based STARKs and
recursive SNARKs based on elliptic curves. STARKs provide post-quantum
security but often with larger proofs or heavier verification \cite{stark}.
Curve recursion can be compact, but its assumptions are not quantum resistant.
Lattice based cryptography offers a third option: post-quantum security from
algebraic hardness assumptions together with commitments that support folding
based recursion \cite{ajtai,latticezkp,latticefold}. Despite recent progress,
there is still no practical, general purpose RISC-V ZkVM built from these
ingredients. The work addresses that gap and asks whether lattice based
folding can support the long computations required for Ethereum block validation.

\section*{Architecture}

The proposed system follows the layered structure common in modern ZkVMs
\cite{sokzkvms}. A guest program is compiled to a 32 bit RISC-V binary and
executed instruction by instruction by an emulator. Each step emits a
trace containing the state of the VM before and after, such as program counter,
register state, decoded instruction, and memory side effects.
This trace is arithmetized into a witness over Goldilocks ring elements.

The core relation is expressed as a Customizable Constraint System (CCS)
\cite{ccs}. A step circuit $F$ enforces correct RISC-V transitions, including
decoding, operand validation, arithmetic or control flow semantics, and
committed memory page updates. To make long executions provable with folding,
the $F$ is wrapped with an augmented circuit $F'$ that instantiates incrementally
verifiable computation (IVC) \cite{valiantivc,nova,hypernova}. For each instruction,
$F'$ checks a commitment to the machine state at the end of last step,
verifies the previous folding proof, and enforces the current RISC-V step. Off
circuit, the prover uses LatticeFold to compress the current step and the
running accumulator into a new accumulator \cite{latticefold}. This replaces
full recursive proof verification at every step with a cheaper folding
operation under lattice assumptions.

Memory is handled with deferred checking rather than a full in circuit RAM
argument. Every load and store is appended to a transcript during execution,
and the final proof verifies permutation, ordering, and read over write
properties using offline memory checking techniques \cite{memorycheck}. To
produce a succinct artifact for an external verifier, the final folded
instance is wrapped with a SuperSpartan style IOP combined with Greyhound
polynomial commitments \cite{spartan,greyhound}. Thanks to Greyhound, final
proofs sizes are in the tens of kilobytes, much smaller than the hundreds of kilobytes
to about one megabyte reported for current ZkVMs.

\section*{Implementation and evaluation}

An incomplete Rust prototype is presented for the proving pipeline.
The implementation is sufficient for initial experiments. The first benchmark
uses a small program that computes the 100th Fibonacci number. This benchmark
was evaluated on the prototype implementation together with Pico\footnote{\url{https://pico-docs.brevis.network/}.},
SP1\footnote{\url{https://docs.succinct.xyz/docs/sp1/introduction}.}, and ZKsync
Airbender\footnote{\url{https://www.zksync.io/airbender}.}
on the same machine. In the prototype implementation, the benchmark
corresponds to 16 executed instructions and 2 memory operations. The benchmark
summary is shown in Table 1.

\begin{center}
	\footnotesize
	\begin{tabular}{|l|c|c|}
		\hline
		\textbf{ZkVM}            & \textbf{Proof time} & \textbf{Peak RAM} \\
		\hline
		Prototype implementation & 11:28.92            & 12.34 GB          \\
		Pico                     & 2:04.50             & 6.27 GB           \\
		SP1                      & 1:50.17             & 9.96 GB           \\
		ZkSync Airbender         & 4:06.73             & 24.34 GB          \\
		\hline
	\end{tabular}

	\smallskip
	\textbf{Table 1.} Fibonacci benchmark comparison.
\end{center}

The results show that the prototype implementation is slower than Pico,
SP1, and ZKsync Airbender on this small benchmark. At the same time, its
memory use stays below Airbender and within a manageable range. The gap in
proving time is attributed to multilinear extension, decomposition, and
folding work in the proving backend rather than to the VM frontend itself.

\paragraph{REVM experiment.} The longer \texttt{revm} based experiment
gives the most encouraging signal for the Ethereum use case. That run
continued for more than 9 hours while the IVC chain kept advancing and memory
usage remained bounded at roughly 13.7 GB. This does not yet prove readiness
for production Ethereum block proving. However, it does show that the design
can sustain a much longer recursive execution without obvious runaway memory
growth or a breakdown of the folding process. In that sense, the experiment is
an encouraging sign that witness norm management and recursive state growth are
not fundamental blockers for the approach.

\paragraph{Scientific questions.} Answers to the research
questions are cautious but positive. The evaluation asks three things: whether
it is feasible to instantiate a full RISC-V execution environment using CCS
compatible with LatticeFold, how a lattice based folding scheme compares in
performance with established ZkVMs, and whether noise growth in lattice
cryptography hinders the ability to prove long traces such as Ethereum block
validation. The first question is answered partially yes, because the
main building blocks can be connected into a single proving pipeline even
though the implementation remains incomplete. The second question currently
receives a negative practical answer, since the prototype is slower than
established ZkVMs. Still, the design retains an important upside: unlike most
current systems, it targets compact post-quantum proofs in the tens of
kilobytes. The third question receives a cautiously positive answer, because
the long \texttt{revm} run suggests that long executions are viable enough to
justify further work, even though production scale Ethereum block proving
remains out of reach.

\paragraph{Identified problems.} The evaluation indicates that the main
cost is not VM execution itself, but multilinear extension, decomposition, and
folding work in the proving backend. The implementation is CPU only and does
not yet benefit from SIMD or GPU acceleration. CCS is expressive, but it does
not appear to be the best model for a VM together with a recursive verifier,
especially for cryptographic components and self referential verifier logic.
CCS is convenient for some linear instruction
semantics, but becomes awkward for parts such as Poseidon2 and the folding
verifier, where the constraint structure starts to work against efficiency.
Taken together, these observations suggest that the current bottlenecks are
mostly engineering and arithmetization issues rather than evidence that
lattice based folding itself is unsuitable.

\paragraph{Future work.} The most direct next steps are to replace
LatticeFold with improved backends such as LatticeFold+ or Neo, revisit CCS as
the computation model for recursion and VM constraints, and exploit SIMD or GPU
acceleration for the regular linear algebra workload
\cite{latticefoldplus,neo}. Beyond that, further work is motivated on better
algebraic layouts for multilinear extensions, reuse of intermediate values
across folding steps, and a proving architecture that exposes more parallelism.
The broader conclusion is therefore not that the approach is already practical,
but that it is concrete enough to optimize and compare as a serious post-
quantum direction for Ethereum execution proving.

\paragraph{Conclusion.} The work presents a feasible path toward Ethereum
proving with lattices. It combines a RISC-V execution model, CCS, folding, and
Greyhound into a design that promises compact proofs, while the current
prototype shows where the remaining performance gaps lie. The main takeaway is
simple: the approach is not ready yet, but it is concrete, measurable, and
worth further optimization.

\newpage

\begin{thebibliography}{99}

	\bibitem{buterin}
	V. Buterin.
	\newblock \emph{Ethereum: A Next-Generation Smart Contract and Decentralized Application Platform}.
	\newblock 2014.

	\bibitem{sokzkvms}
	Y. Yang, Y. Cheng, H. Tang, G. Yang, B. Zhang, and K. Ren.
	\newblock \emph{SoK: Understanding zkVM: From Research to Practice}.
	\newblock Cryptology ePrint Archive, Paper 2026/525, 2026.

	\bibitem{stark}
	E. Ben-Sasson, I. Bentov, Y. Horesh, and M. Riabzev.
	\newblock \emph{Scalable, Transparent, and Post-Quantum Secure Computational Integrity}.
	\newblock Cryptology ePrint Archive, Paper 2018/046, 2018.

	\bibitem{ajtai}
	M. Ajtai.
	\newblock \emph{Generating Hard Instances of Lattice Problems (Extended Abstract)}.
	\newblock In \emph{Proceedings of STOC}, pages 99--108, 1996.

	\bibitem{latticezkp}
	V. Lyubashevsky, N. K. Nguyen, and M. Plan\c{c}on.
	\newblock \emph{Lattice-Based Zero-Knowledge Proofs and Applications: Shorter, Simpler, and More General}.
	\newblock Cryptology ePrint Archive, Paper 2022/284, 2022.

	\bibitem{latticefold}
	D. Boneh and B. Chen.
	\newblock \emph{LatticeFold: A Lattice-based Folding Scheme and its Applications to Succinct Proof Systems}.
	\newblock Cryptology ePrint Archive, Paper 2024/257, 2024.

	\bibitem{ccs}
	S. Setty, J. Thaler, and R. Wahby.
	\newblock \emph{Customizable Constraint Systems for Succinct Arguments}.
	\newblock Cryptology ePrint Archive, Paper 2023/552, 2023.

	\bibitem{valiantivc}
	P. Valiant.
	\newblock \emph{Incrementally Verifiable Computation or Proofs of Knowledge Imply Time/Space Efficiency}.
	\newblock In \emph{Theory of Cryptography}, pages 1--18, 2008.

	\bibitem{nova}
	A. Kothapalli, S. Setty, and I. Tzialla.
	\newblock \emph{Nova: Recursive Zero-Knowledge Arguments from Folding Schemes}.
	\newblock Cryptology ePrint Archive, Paper 2021/370, 2021.

	\bibitem{hypernova}
	A. Kothapalli and S. Setty.
	\newblock \emph{HyperNova: Recursive Arguments for Customizable Constraint Systems}.
	\newblock Cryptology ePrint Archive, Paper 2023/573, 2023.

	\bibitem{memorycheck}
	A. Ozdemir, E. Laufer, and D. Boneh.
	\newblock \emph{Volatile and Persistent Memory for zkSNARKs via Algebraic Interactive Proofs}.
	\newblock Cryptology ePrint Archive, Paper 2024/979, 2024.

	\bibitem{spartan}
	S. Setty.
	\newblock \emph{Spartan: Efficient and General-Purpose zkSNARKs without Trusted Setup}.
	\newblock Cryptology ePrint Archive, Paper 2019/550, 2019.

	\bibitem{greyhound}
	N. K. Nguyen and G. Seiler.
	\newblock \emph{Greyhound: Fast Polynomial Commitments from Lattices}.
	\newblock Cryptology ePrint Archive, Paper 2024/1293, 2024.

	\bibitem{latticefoldplus}
	D. Boneh and B. Chen.
	\newblock \emph{LatticeFold+: Faster, Simpler, Shorter Lattice-Based Folding for Succinct Proof Systems}.
	\newblock Cryptology ePrint Archive, Paper 2025/247, 2025.

	\bibitem{neo}
	W. Nguyen and S. Setty.
	\newblock \emph{Neo: Lattice-Based Folding Scheme for CCS over Small Fields and Pay-Per-Bit Commitments}.
	\newblock Cryptology ePrint Archive, Paper 2025/294, 2025.

\end{thebibliography}

\end{document}
